\documentclass{article} % For LaTeX2e
\usepackage[T1]{fontenc}
\usepackage{iclr2027_conference,times}

% Optional math commands from https://github.com/goodfeli/dlbook_notation.
\input{math_commands.tex}

\usepackage{hyperref}
\usepackage{url}
\usepackage{amsthm}
\usepackage{enumitem}
\usepackage{booktabs}
\usepackage{graphicx}

% theorem environments (as in Morwani et al., 2024)
\newtheorem{theorem}{Theorem}
\newtheorem*{theorem*}{Theorem}
\newtheorem{lemma}[theorem]{Lemma}
\newtheorem*{lemma*}{Lemma}
\newtheorem{proposition}[theorem]{Proposition}
\newtheorem{corollary}[theorem]{Corollary}
\newtheorem{conjecture}[theorem]{Conjecture}
\newtheorem*{informalthm}{Theorem (informal)}
\theoremstyle{definition}
\newtheorem{definition}[theorem]{Definition}
\newtheorem{assumption}[theorem]{Assumption}
\theoremstyle{remark}
\newtheorem{remark}[theorem]{Remark}
\newtheorem*{remark*}{Remark}

% notation
\newcommand{\Fp}{F}                                % positive frequencies {1,...,(p-1)/2}
\newcommand{\spt}{\mathrm{spt}}
\newcommand{\cov}[2]{\kappa_{#1}(#2)}              % coverage of frequency #1 at offset #2
\newcommand{\Kstar}{K^{*}}
\newcommand{\curlyx}{\mathcal{X}}
\newcommand{\curlyy}{\mathcal{Y}}
\newcommand{\curlyp}{\mathcal{P}}
\newcommand{\lnorm}[2]{\Vert #1 \Vert_{#2}}

\title{Explaining Feature Sparsity Through the\\ Margin--Regularization Trade-off}

% Authors must not appear in the submitted version.
\author{Anonymous authors}

\newcommand{\fix}{\marginpar{FIX}}
\newcommand{\new}{\marginpar{NEW}}

%\iclrfinalcopy % Uncomment for camera-ready version, but NOT for submission.
\begin{document}

\maketitle

\begin{abstract}
Explaining the number of features in a network has been a long-desired result in mechanistic interpretability. \citet{morwani2024feature} laid the baseline with their principle of margin-maximization, showing provably that in networks where the regularization ($L_{2,\nu}$) matches its homogeneity degree ($\nu$), all features would be used for modular addition or group composition tasks. Empirically, however, networks trained with ordinary weight decay consistently show a sparse selection of features. This paper investigates why.
Our inquiry focuses specifically on modular addition and group composition tasks in polynomial and multilinear models. Our primary theoretical finding centers on the cost of new features---a trade-off between the logit margin of the solution and regularization. While more features improve margin, they cost the model in regularization. Hence we show that an optimal number of features can be calculated for networks based on the task, their homogeneity degree and the regularization. Further, our theory explains the transition in the number of features from dense to sparse---as seen empirically in \citet{mccracken2025}---by showing that this transition arises on changing either the homogeneity degree or regularization of the network.
Specifically, we empirically show that the normalized margin attained by these networks peaks at the predicted optimum for the network, matching the result of \citet{morwani2024feature} for their specific case. The optimums are also consistent with the empirical results seen in \citet{chughtai2023toy}.
\end{abstract}

\input{sections/intro}
\input{sections/prelim}
\input{sections/approach}
\input{sections/cyclic}
% \input{sections/groups} 
\input{sections/experiments}
\input{sections/discussion}

% Required for ICLR 2027; author will supply the AI-use statement before submission.
% Place it here as an unnumbered subsection, before the references (one page maximum).

\subsection*{Reproducibility statement}
All theorems are stated with their hypotheses in Sections~\ref{sec:approach} and~\ref{sec:cyclic} and proved in full in the Appendix; the constants of every construction were checked numerically to machine precision. Experimental settings are listed in Section~\ref{sec:experiments}; code and configuration files are provided as supplementary material.

\clearpage
\bibliography{refs}
\bibliographystyle{iclr2027_conference}

\newpage
\appendix
\input{sections/appendix}
\input{sections/appendix_experiments}

\end{document}
